\section{Convergence under $H_A$}
\begin{frame}{\hspace{3cm}}
\centering

\vspace{1.5cm}

\begin{beamercolorbox}[sep=20pt,center,rounded=true,shadow=true]{title}
    {\fontsize{20}{24}\selectfont \textbf{Convergence under $H_A$}}
\end{beamercolorbox}

\vspace{0.8cm}


\vfill

\end{frame}
%------------------------------------------------

\begin{frame}{Setup under $H_A$}
\begin{itemize}
    \item True model is nonparametric:
    \[
    f(x) \neq f(x,\theta)
    \]

    \item Goal:
    \begin{itemize}
        \item Show estimator behaves like kernel estimator
    \end{itemize}
\end{itemize}
\end{frame}

%------------------------------------------------

\begin{frame}{Theorem under $H_A$}
\small
\textbf{Theorem 2.} Under the assumptions stated previously
\[
\left|
\frac{f(x_i,\theta_A) - f(x_i,\hat{\theta})}{f(x_i,\theta_A)}
\right|
\leq
|\hat{\theta} - \theta_A| R(x_i),
\]
where $E[R(X_i)]^2 < \infty$. Then
\[
a_n^{-1} \int |g(x,\hat{\pi}) - f(x)| \, dx = O_p(1).
\tag{d}
\]
\end{frame}

%------------------------------------------------

\begin{frame}{Proof Idea}
\small
\textbf{Proof.} To prove (d), we show that
\[
\frac{1}{n} \sum \log g(x_i,\pi)
=
- O\left(
\int \left|\sqrt{g_A(x,\pi)} - \sqrt{f(x)}\right|^2 dx
\right)
+ O_p(a_n^2),
\tag{e}
\]
where
\[
g_A(x,\pi) = \pi f(x,\theta_A) + (1-\pi) f(x).
\]

\vspace{0.3cm}

\begin{itemize}
    \item Result (e) is stronger than (d)
    \item $\hat{\pi}$ maximizes (e)
    \item Equivalent to minimizing Hellinger distance
\end{itemize}
\end{frame}
%------------------------------------------------

\begin{frame}{Step 1}
\begin{itemize}
    \item Show:
    \[
    \frac{1}{n}\sum \log g(x_i,\pi)
    =
    \mathbb{E}[\log g_A(X,\pi)] + O_p(n^{-1/2})
    \]

    \item Method:
    \begin{itemize}
        \item Taylor expansion around $\theta_A$
    \end{itemize}
\end{itemize}
\end{frame}

%------------------------------------------------

\begin{frame}{Step 2}
\begin{itemize}
    \item Show:
    \[
    \mathbb{E}[\log g_A(X,\pi)]
    =
    -O\Big(\int (\sqrt{g(x,\pi)}-\sqrt{f(x)})^2 dx\Big)
    \]

    \item Interpretation:
    \begin{itemize}
        \item Hellinger distance
    \end{itemize}

\textbf{Step 1} and \textbf{Step 2} together imply that 
\[
\hat{\pi} = O_p(n^{-1/4}).
\]
By iterating this process we ultimately achieve (e).

\end{itemize}
\end{frame}

%------------------------------------------------

\begin{frame}{Final Result}
\begin{itemize}
    \item Using the expansion of $
H(\pi) = \frac{1}{\sqrt{n}} \sum \left[ \log g_A(x_i,\pi) - E \log g_A(X,\pi) \right].
$ about $\pi = 0$ and using an iterating argument from previous calculations(from which we know that $\hat\pi = O_p(max(a_n,n^{-1/2}))$), it can be shown $\hat\pi = O_p(a_n)$ provided that $a_n n^{\alpha/2} \to \infty $ for some $0<\alpha<1$.

    \item Hence, Non-parametric model dominates.
    
\end{itemize}
\end{frame}
